\documentclass[a4paper,12pt]{article}
\usepackage[bibliografia]{plantilla}

\titulo{Guía rápida de la plantilla}
\autor{Carlos Sesma}
\subtitulo{Cómo escribir documentos en este repositorio}

\begin{document}

\portada
\indice

\section{Qué hay aquí}

Este repositorio aloja varios documentos LaTeX que comparten estilo. Cada uno
vive en su propia carpeta dentro de \texttt{documentos/}, con su bibliografía y
sus imágenes, y todos los PDF se generan en \texttt{out/}.

El preámbulo común está en \texttt{comun/plantilla.sty}, así que un documento
no repite la lista de paquetes: se limita a cargar la plantilla y a declarar
sus datos.

\section{Crear un documento}

\begin{lstlisting}[language=bash,caption={Crear y compilar}]
./doc nuevo       # pregunta plantilla, titulo y autor
./doc compila     # genera los PDF en out/
./doc lista       # que hay y que esta compilado
\end{lstlisting}

\begin{nota}
    \texttt{./doc nuevo} también admite argumentos
    (\texttt{-{}-plantilla}, \texttt{-{}-titulo}, \texttt{-{}-autor}) para
    usarlo desde un script.
\end{nota}

\section{Plantillas}

Las plantillas son carpetas en \texttt{plantillas/}. Este repositorio añade una
propia, \texttt{receta}, que no venía en la original: basta con crear la carpeta
con un \texttt{documento.tex} y un \texttt{plantilla.conf}, y la CLI la ofrece
en el menú sin tocar una línea de código.

\section{Qué aporta la plantilla}

Además de los paquetes, \texttt{plantilla.sty} define la portada, el índice y
las cajas destacadas.

\begin{aviso}
    Los documentos que no citan nada se declaran con
    \texttt{\textbackslash usepackage\{plantilla\}} a secas, sin la opción
    \texttt{bibliografia}, y así no arrastran biber.
\end{aviso}

Para el detalle tipográfico, la referencia sigue siendo la de siempre
\cite{knuth1984}.

\printbibliography

\end{document}
